\section{Conclusions}\label{sec:conclusion}
This year-long, multi-seed benchmark of residential cooling control for a Pune
apartment supports five conclusions.
\begin{enumerate}
\item With a realistic thermostat as the baseline and a fair discomfort metric,
      advanced controllers pay off mainly in comfort and fewer compressor starts, and
      (except the DQN) in lower peak demand. They do not reduce cost relative to the
      thermostat.
\item Among controllers that deliver good comfort, the hierarchical Hybrid (economic MPC
      with a PI correction layer and safety override) is best. It costs
      \RelHybPidCost\,\% less than SIMC-tuned PID and \RelHybMpcCost\,\% less than
      E-MPC, with \RelHybPidKh\,\% and \RelHybMpcKh\,\% less discomfort respectively.
      The ranking holds under a +3\,K heatwave, thermal mass, the tested tariffs (on comfort) and
      a \UqN-sample variation of eleven assumed inputs. The size of the saving over
      PID varies with the inputs ($\UqHPcostLo$\ to $\UqHPcostHi$\,\%). The PI layer
      provides the cost saving, and the safety override provides most of the
      comfort gain.
\item A DQN with a safety shield, multi-seed training and a separate validation set is
      competitive. Without the shield it is not, and its performance varies with the
      training seed.
\item A hypothetical three-tier ToD tariff with a 1.7:1 price ratio produces little load
      shifting, because peak-window cooling serves occupied zones and coincides with
      heat release from the roof slab. A 3.6:1 ratio shifts about 5 percentage points of
      cooling energy. Under the actual MSEDCL residential structure the Hybrid
      remains the most comfortable and 4\,\% cheaper than PID, with the DQN within
      0.1\,\% of it on cost.
\item Indoor air quality depends on ventilation and filtration design, not on the
      thermal controller. Recirculating HEPA filtration keeps PM$_{2.5}$ well below
      outdoor levels. A demand-controlled dehumidifying DOAS removes most humid hours,
      at an energy cost of about a quarter of annual electricity that deserves dedicated
      control.
\end{enumerate}
Future work should calibrate the model against measured data or EnergyPlus, add
measured hourly PM$_{2.5}$, use adaptive and humidity-aware comfort metrics, and
co-optimise ventilation with cooling in one MPC.
